%%%%%%%%%%%%%%%%%%%%%%%%%%%%%%%%%%%%

\section{少ない母集団からの標本抽出}

%%%%%%%%%%%%%%%%%%%%%%%%%%%%%%%%%%%%

\begin{frame}[shrink=10]
\frametitle{復元抽出}

\hl{復元抽出}では、引いたものを元に戻す。

\begin{itemize}

\item 赤チップ5個、青チップ3個、オレンジチップ2個が入った袋があるとする。最初に引いたチップが青である確率は何か？
\begin{center}
5 \textcolor{red}{$\CIRCLE$}~, 3 \textcolor{blue}{$\CIRCLE$}~, 2 \textcolor{orange}{$\CIRCLE$}
\end{center}

\[ Prob(1\text{番目のチップが青}) = \frac{3}{5 + 3 + 2} = \frac{3}{10} = 0.3 \]

\item 1回目の引きで青チップを引いたとする。復元抽出で引く場合、2回目の引きで青チップを引く確率は何か？

\begin{center}
$1\text{回目}$: 5 \textcolor{red}{$\CIRCLE$}~, 3 \textcolor{blue}{$\CIRCLE$}~, 2 \textcolor{orange}{$\CIRCLE$} \\

$2\text{回目}$: 5 \textcolor{red}{$\CIRCLE$}~, 3 \textcolor{blue}{$\CIRCLE$}~, 2 \textcolor{orange}{$\CIRCLE$}
\end{center}

\[ Prob(2\text{番目のチップが青} | 1\text{番目のチップが青}) = \frac{3}{10} = 0.3 \]

\end{itemize}

\end{frame}

%%%%%%%%%%%%%%%%%%%%%%%%%%%%%%%%%%%%%

\begin{frame}
\frametitle{復元抽出（続き）}

\begin{itemize}

\item 1回目の引きでオレンジチップを引いたとする。復元抽出で引く場合、2回目の引きで青チップを引く確率は何か？

\begin{center}
$1\text{回目}$: 5 \textcolor{red}{$\CIRCLE$}~, 3 \textcolor{blue}{$\CIRCLE$}~, 2 \textcolor{orange}{$\CIRCLE$} \\
$2\text{回目}$: 5 \textcolor{red}{$\CIRCLE$}~, 3 \textcolor{blue}{$\CIRCLE$}~, 2 \textcolor{orange}{$\CIRCLE$}
\end{center}
\[ Prob(2\text{番目のチップが青} | 1\text{番目のチップがオレンジ}) = \frac{3}{10} = 0.3 \]

\item 復元抽出で、2回続けて青チップを引く確率は何か？
\begin{center}

$1\text{回目}$: 5 \textcolor{red}{$\CIRCLE$}~, 3 \textcolor{blue}{$\CIRCLE$}~, 2 \textcolor{orange}{$\CIRCLE$} \\
$2\text{回目}$: 5 \textcolor{red}{$\CIRCLE$}~, 3 \textcolor{blue}{$\CIRCLE$}~, 2 \textcolor{orange}{$\CIRCLE$}
\end{center}
\footnotesize
\[ Prob(1\text{番目のチップが青}) \cdot Prob(2\text{番目のチップが青} | 1\text{番目のチップが青}) \]
\[ = 0.3 \times 0.3 \]
\[ = 0.3^2 = 0.09 \]

\end{itemize}

\end{frame}

%%%%%%%%%%%%%%%%%%%%%%%%%%%%%%%%%%%%%

\begin{frame}[shrink=20]
\frametitle{復元抽出（続き）}

\begin{itemize}

\item 復元抽出では、2番目のチップが青である確率は、1番目のチップの色に依存しない。なぜなら、1回目に引いたものを袋に戻すからである。
\[ Prob(\text{青} | \text{青}) = Prob(\text{青} | \text{オレンジ}) \]

\item さらに、この確率は1番目のチップが青である確率と等しい。なぜなら、復元抽出では袋の構成が変わらないからである。
\[ Prob(\text{青} | \text{青}) = Prob(\text{青}) \]

\item \hl{復元抽出では、引き操作は独立である。}

\end{itemize}

\end{frame}


%%%%%%%%%%%%%%%%%%%%%%%%%%%%%%%%%%%%%

\begin{frame}
\frametitle{非復元抽出}

\hl{非復元抽出}では、引いたものを元に戻さない。

\begin{itemize}

\item 1回目の引きで青チップを引いたとする。非復元抽出で引く場合、2回目の引きで青チップを引く確率は何か？
\begin{center}
$1\text{回目}$: 5 \textcolor{red}{$\CIRCLE$}~, 3 \textcolor{blue}{$\CIRCLE$}~, 2 \textcolor{orange}{$\CIRCLE$} \\
$2\text{回目}$: 5 \textcolor{red}{$\CIRCLE$}~, 2 \textcolor{blue}{$\CIRCLE$}~, 2 \textcolor{orange}{$\CIRCLE$}
\end{center}
\[ Prob(2\text{番目のチップが青} | 1\text{番目のチップが青}) = \frac{2}{9} = 0.22 \]

\item 非復元抽出で、2回続けて青チップを引く確率は何か？
\begin{center}

$1\text{回目}$: 5 \textcolor{red}{$\CIRCLE$}~, 3 \textcolor{blue}{$\CIRCLE$}~, 2 \textcolor{orange}{$\CIRCLE$} \\
$2\text{回目}$: 5 \textcolor{red}{$\CIRCLE$}~, 2 \textcolor{blue}{$\CIRCLE$}~, 2 \textcolor{orange}{$\CIRCLE$}
\end{center}
\footnotesize
\[ Prob(1\text{番目のチップが青}) \cdot Prob(2\text{番目のチップが青} | 1\text{番目のチップが青}) \]
\[  = 0.3 \times 0.22 \]
\[ = 0.066 \]

\end{itemize}

\end{frame}

%%%%%%%%%%%%%%%%%%%%%%%%%%%%%%%%%%%%%

\begin{frame}
\frametitle{非復元抽出（続き）}

\begin{itemize}

\item 非復元抽出では、1番目が青であった場合の2番目のチップが青である確率は、1番目のチップが青である確率と等しくない。なぜなら、袋の構成が1回目の引きの結果によって変わるからである。
\[ Prob(\text{青} | \text{青}) \ne Prob(\text{青}) \]

\item \hl{非復元抽出では、引き操作は独立でない。}

\item これは、標本サイズが小さいときに特に注意が必要である。たとえば、（巨大な）袋に10,000個のチップが入っている場合、1つのチップを引き出しても2回目の引きの確率にはそれほど大きな影響を与えない。

\end{itemize}

\end{frame}

%%%%%%%%%%%%%%%%%%%%%%%%%%%%%%%%%%%%%

\begin{frame}
\frametitle{練習問題}

\pq{ほとんどのカードゲームでは、カードは非復元で配られる。エースを引いてから3を引く確率は何か？最も近い答えを選びなさい。}

\twocol{0.3}{0.6}{
\begin{enumerate}[(a)]
\item 0.0045
\item 0.0059
\solnMult{0.0060}
\item 0.1553
\end{enumerate}
}
{
\soln{
\[ P(\text{エース次に3}) = \frac{4}{52} \times \frac{4}{51} \approx 0.0060 \]
}
}

\end{frame}

%%%%%%%%%%%%%%%%%%%%%%%%%%%%%%%%%%%%%
